\chapter{Vaginal Odor}

\section*{Definition}
Vaginal odor varies with menstruation, sweat, sex, and pH. A persistent strong or fishy odor, especially with discharge, itching, or burning, may indicate infection or a retained foreign body.

\section*{When to See a Doctor}
Seek urgent evaluation for Fever, pelvic pain, pregnancy, bleeding, severe pain, or foul odor with systemic illness requires prompt care. Arrange routine or prompt assessment for persistent, recurrent, unexplained, or function-limiting symptoms.

\section*{How It Develops}
Vaginal odor can result from normal variation, bacterial vaginosis, trichomoniasis, retained foreign material, a sexually transmitted infection, or less commonly a fistula or malignancy. Discharge, itching, pain, bleeding, pregnancy, and sexual exposure guide evaluation.

\section*{Causes}
\begin{itemize}
  \item Bacterial vaginosis, trichomoniasis, other vaginitis, retained tampon, rectovaginal fistula, and rarely vaginal or cervical cancer.
  \item Medication effects, environmental exposures, and chronic disease may contribute.
  \item Less common but important causes include malignancy, vascular disease, or organ failure when symptoms persist or progress.
\end{itemize}

\section*{Related Symptoms}
\begin{itemize}
  \item Pain, fever, fatigue, nausea, or changes in normal organ function
  \item Bleeding, weight change, shortness of breath, weakness, or neurologic symptoms when a serious cause is present
\end{itemize}

\section*{Medical Evaluation}
\begin{itemize}
  \item History addresses onset, duration, severity, triggers, exposures, medications, medical conditions, pregnancy status when relevant, and warning symptoms.
  \item Examination focuses on vital signs and the organ systems suggested by the symptom.
  \item Testing may include blood or urine studies, cultures, imaging, physiologic testing, or specialist examination according to the clinical pattern.
\end{itemize}

\section*{Treatment Options}
Treatment is based on examination and appropriate testing. Confirmed bacterial vaginosis, trichomoniasis, or another infection receives recommended therapy; avoid douching and do not self-treat recurrent odor with leftover antibiotics.

\section*{Self-Care and Home Management}
Avoid known triggers, maintain reasonable hydration, record timing and associated symptoms, and use only treatments appropriate for the suspected cause. Do not use leftover antibiotics or delay urgent care because symptoms temporarily improve.

\section*{References}
\begin{thebibliography}{99}
\bibitem{vaginal_odor:cdc} Workowski KA, Bachmann LH, Chan PA, et al. Sexually Transmitted Infections Treatment Guidelines, 2021. \textit{MMWR Recomm Rep}. 2021;70(4):1--187.
\bibitem{vaginal_odor:acog} American College of Obstetricians and Gynecologists. Vaginitis in Nonpregnant Patients. ACOG Practice Bulletin No. 215. \textit{Obstet Gynecol}. 2020;135:e1--e17.
\end{thebibliography}
